% GENERATED FILE. Edit canonical lessons, then run npm run generate:books.
% canonical-source-hash: fnv1a64:bac7f95965aee781
% canonical-lessons: PA-W10-form-order, PA-W10-form-select, PA-W10-form-supported, PA-W10-form-delayed, PA-W10-form-repair-values, PA-W10-form-repair-layout, PA-W10-form-no-model

\chapter{Join Three Fictional Form Fields}
\label{ch:pa-form-age-phone-date-checkpoint}
\hlchaptermodality{154}

\begin{chapteropening}
\textbf{By the end of this chapter:} \emph{I can choose taught fictional age, phone, and date values from one nonverbal cue, write the three Gurmukhi fields without a visible answer model, and check each error dimension separately.}

The learner fills age, phone, and date lines from one remembered selector without romanization or copyable answers, then repairs selection, digit order, date format, spelling, spacing, and placement in separate passes.
\end{chapteropening}

\section[\texorpdfstring{\pa{ਉਮਰ}: · \pa{ਫ਼ੋਨ} … (three familiar …)}{ਉਮਰ: · ਫ਼ੋਨ … (three familiar …)}]{Three known labels, one gentle order}

\label{lesson:PA-W10-form-order}



\begin{quote}
Point to the three labels you have already written: \textbf{\pa{ਉਮਰ}:}, \textbf{\pa{ਫ਼ੋਨ}:}, \textbf{\pa{ਤਾਰੀਖ਼}:}. No new field or value is added here.
\end{quote}

\subsection*{Writing — first labels, then values}
This little fictional form puts age on the first line, phone on the second, and date on the third. Copy only the three labels in that order, with one empty space after each colon. Leave the values blank for now.

\begin{quote}
\pa{ਉਮਰ}: \_\_\_\_\_\_
\end{quote}

\begin{quote}
\pa{ਫ਼ੋਨ}: \_\_\_\_\_\_
\end{quote}

\begin{quote}
\pa{ਤਾਰੀਖ਼}: \_\_\_\_\_\_
\end{quote}

\subsection*{Your turn — cover and point}
Cover the form. Point to three empty lines and say which known field belongs on each. Uncover it to check only the line order, not any digits.

\begin{tcolorbox}[breakable,colback=teal!4,colframe=teal!35!black,title={Before you move on}]
Today only the labels and their places matter. The next step reunites the small closed banks already learned one at a time.
\end{tcolorbox}
\section[\texorpdfstring{\pa{ਕ} / \pa{ਖ} (two fictional cards)}{ਕ / ਖ (two fictional cards)}]{One card chooses three already-known values}

\label{lesson:PA-W10-form-select}



\begin{quote}
The old closed banks are visible again. Age A is \textbf{\pa{੧੫}}, B is \textbf{\pa{੨੫}}. Phone A is \textbf{\pa{੦੧੨} \pa{੨੫੧}}, B is \textbf{\pa{੦੨੫} \pa{੧੨੫}}. Date A is \textbf{\pa{੧੫}/\pa{੦੧}/\pa{੨੦੨੫}}, B is \textbf{\pa{੨੫}/\pa{੦੨}/\pa{੨੦੨੫}}. These are fictional practice values, not personal data.
\end{quote}

\subsection*{You'll want to know — a single shared cue}
On this practice card \textbf{\pa{ਕ}} chooses A from \emph{each} bank; \textbf{\pa{ਖ}} chooses B from \emph{each} bank. You already met these shapes in phone and date. For age, the card now points to its old A or B bank too. The card mark is never copied into a form field.

\subsection*{Your turn — select before writing}
Point to \textbf{\pa{ਕ}}. Point to the A value in each bank, top to bottom: age, phone, date. Do not write yet. Repeat for \textbf{\pa{ਖ}} and its B values. Check selection before checking any digits.

\begin{tcolorbox}[breakable,colback=teal!4,colframe=teal!35!black,title={Before you move on}]
The banks stay visible here. Selection is the only new step.
\end{tcolorbox}
\section[\texorpdfstring{\pa{ਉਮਰ}: \pa{੧੫} · \pa{ਫ਼ੋਨ} … (supported …)}{ਉਮਰ: ੧੫ · ਫ਼ੋਨ … (supported …)}]{Fill the little form with the bank still open}

\label{lesson:PA-W10-form-supported}



\begin{quote}
The card \textbf{\pa{ਕ}} asks for A in each old bank. The three A values are visible: age \textbf{\pa{੧੫}}, phone \textbf{\pa{੦੧੨} \pa{੨੫੧}}, date \textbf{\pa{੧੫}/\pa{੦੧}/\pa{੨੦੨੫}}.
\end{quote}

\subsection*{Writing — one familiar line at a time}
For \textbf{\pa{ਕ}}, copy the selected value after each printed label. Finish one line before moving to the next. Keep the phone's single group space and the date's two slash marks.

\begin{quote}
\pa{ਉਮਰ}: \_\_\_\_\_\_
\end{quote}

\begin{quote}
\pa{ਫ਼ੋਨ}: \_\_\_\_\_\_
\end{quote}

\begin{quote}
\pa{ਤਾਰੀਖ਼}: \_\_\_\_\_\_
\end{quote}

\subsection*{Your turn — three small checks}
Compare the selected bank letter first, the value on each line second, and line order third. Correct just one difference at a time. This supported copy does not count as independent writing.

\begin{tcolorbox}[breakable,colback=teal!4,colframe=teal!35!black,title={Before you move on}]
The bank remains visible throughout this gentle supported attempt.
\end{tcolorbox}
\section[\texorpdfstring{\pa{ਉਮਰ}: · \pa{ਫ਼ੋਨ} … (delayed …)}{ਉਮਰ: · ਫ਼ੋਨ … (delayed …)}]{Hide the bank for one short turn}

\label{lesson:PA-W10-form-delayed}



\begin{quote}
Read the familiar B set once: \textbf{\pa{ਉਮਰ}: \pa{੨੫}}, \textbf{\pa{ਫ਼ੋਨ}: \pa{੦੨੫} \pa{੧੨੫}}, \textbf{\pa{ਤਾਰੀਖ਼}: \pa{੨੫}/\pa{੦੨}/\pa{੨੦੨੫}}. Point to one line at a time.
\end{quote}

\subsection*{Writing — brief delayed copy}
Cover the model. Wait five seconds. The card is \textbf{\pa{ਖ}}. Write the three lines in the known order, recalling just one line at a time. Keep the model covered until all three lines are finished.

\begin{quote}
\pa{ਉਮਰ}: \_\_\_\_\_\_
\end{quote}

\begin{quote}
\pa{ਫ਼ੋਨ}: \_\_\_\_\_\_
\end{quote}

\begin{quote}
\pa{ਤਾਰੀਖ਼}: \_\_\_\_\_\_
\end{quote}

\subsection*{Your turn — uncover only to compare}
Uncover the model. Compare age first, phone second, and date third. This is still delayed-copy practice, not an independent checkpoint.

\begin{tcolorbox}[breakable,colback=teal!4,colframe=teal!35!black,title={Before you move on}]
If one line was hard, revisit only that old field before continuing.
\end{tcolorbox}
\section[\texorpdfstring{\pa{੨੫} · \pa{੦੨੫} \pa{੧੨੫} … (three value checks)}{੨੫ · ੦੨੫ ੧੨੫ … (three value checks)}]{Check the values in three little passes}

\label{lesson:PA-W10-form-repair-values}



\begin{quote}
The card \textbf{\pa{ਖ}} still selects B for all three fictional fields. Before checking any digit, ask: did every line use B, not a mixture of A and B?
\end{quote}

\begin{grammarlens}[title={Grammar lens — selection, order, format}]
Make three separate passes. First check the selected bank letter for each line. Next check digits left to right, including the phone's group of three and three. Last check the date's day/month/year boxes and two slash marks. Do not erase a correct line to fix another one.
\end{grammarlens}

\subsection*{Your turn — one change at a time}
For \textbf{\pa{ਖ}}, suppose the age line says \textbf{\pa{੧੫}}. Change only its selected value to \textbf{\pa{੨੫}}. If the phone line says \textbf{\pa{੦੨੫} \pa{੧੫੨}}, keep the first group and repair only the last two digit positions. If the date says \textbf{\pa{੨੫}/\pa{੦੨੨੦੨੫}}, keep all digits and add only the missing slash.

\begin{tcolorbox}[breakable,colback=teal!4,colframe=teal!35!black,title={Before you move on}]
Label spelling, spacing, and placement get their own pass next.
\end{tcolorbox}
\section[\texorpdfstring{\pa{ਉਮਰ}: · \pa{ਫ਼ੋਨ} … (label and …)}{ਉਮਰ: · ਫ਼ੋਨ … (label and …)}]{Check the labels and their places}

\label{lesson:PA-W10-form-repair-layout}



\begin{quote}
Read the known labels once: \textbf{\pa{ਉਮਰ}:}, \textbf{\pa{ਫ਼ੋਨ}:}, \textbf{\pa{ਤਾਰੀਖ਼}:}. The little dot marks in \textbf{\pa{ਫ਼}} and \textbf{\pa{ਖ਼}} belong to their spellings.
\end{quote}

\subsection*{Writing — three different layout passes}
First compare label spelling; restore a missing dot without touching digits. Then compare spacing: one readable space follows each colon, and the phone has one group space. Finally compare placement: age is line one, phone line two, date line three. Move a whole correct value to its proper line; do not rewrite its digits.

\subsection*{Your turn — a single spelling repair}
In \textbf{\pa{ਤਾਰੀਖ}: \pa{੨੫}/\pa{੦੨}/\pa{੨੦੨੫}}, fix only the label to \textbf{\pa{ਤਾਰੀਖ਼}}. The date stays unchanged. In \textbf{\pa{ਫੋਨ}: \pa{੦੨੫} \pa{੧੨੫}}, fix only the label to \textbf{\pa{ਫ਼ੋਨ}}.

\begin{tcolorbox}[breakable,colback=teal!4,colframe=teal!35!black,title={Before you move on}]
You have now practised six distinct checks: choice, digit order, date format, label spelling, spacing, and placement. The final short attempt hides the value bank again.
\end{tcolorbox}
\section[\texorpdfstring{\pa{ਉਮਰ}: · \pa{ਫ਼ੋਨ} … (independent …)}{ਉਮਰ: · ਫ਼ੋਨ … (independent …)}]{One no-model fictional three-field form}

\label{lesson:PA-W10-form-no-model}



\begin{quote}
Close the earlier lessons. No value bank, support-language label, Latin-digit version, romanization, or copyable Gurmukhi value appears below.
\end{quote}

\subsection*{Writing — independent controlled choice}
Complete the fictional form for the one card mark. Use only the old known values. Write all three in Gurmukhi digits from memory. Do not reopen a lesson until all three lines are filled.

\begin{quote}
\pa{ਖ}
\end{quote}

\begin{quote}
\pa{ਉਮਰ}: \_\_\_\_\_\_
\end{quote}

\begin{quote}
\pa{ਫ਼ੋਨ}: \_\_\_\_\_\_
\end{quote}

\begin{quote}
\pa{ਤਾਰੀਖ਼}: \_\_\_\_\_\_
\end{quote}

\subsection*{Your turn — repair after the attempt}
Only after writing, check six dimensions separately: selection, digit order, date format, label spelling, spacing, then field placement. Repair the first different dimension without erasing a correct one.

\begin{tcolorbox}[breakable,colback=teal!4,colframe=teal!35!black,title={Before you move on}]
This closes only an untimed three-field checkpoint. The complete six-field form, timed writing, two full mocks, and human A1 validation remain separate dependency-linked work.
\end{tcolorbox}
